\scanpage{72}
% archive sheet number: 73
Soit

\[
\begin{aligned}
u&=\sum\rho_i\bar b_i\otimes a_i & \rho_i&=\rho_i(u)\quad\text{etc}\\
v&=\sum\sigma_j\bar d_j\otimes c_j & \sigma_j&=\sigma_j(v)
\end{aligned}
\]

d'où

\[
\begin{aligned}
(\overset{p}{\Lambda}u)\wedge(\overset{q}{\Lambda}v)
&=\frac1{(p+q)!}\sum_{\substack{i_1\ldots i_p\\j_1\ldots j_q}}\rho_{i_1}\ldots\rho_{i_p}\sigma_{j_1}\ldots\sigma_{j_q}
\overline{b_{i_1}\wedge\ldots b_{i_p}\wedge d_{j_1}\wedge\ldots d_{j_q}}\otimes\\[-2pt]
&\hspace{25mm}(a_{i_1}\wedge\ldots a_{i_p}\wedge c_{j_1}\ldots c_{j_q})
\end{aligned}
\]

On peut sommer sur des ($i_\alpha$) tous distincts et des ($j_\beta$) tous distincts, soit

donc

\[
\begin{aligned}
\lVert(\overset{p}{\Lambda}u)\wedge(\overset{q}{\Lambda}v)\rVert_1
&\leq\frac{1}{(p+q)!}p!q!\sum_{\substack{i_1<\ldots<i_p\\j_1<\ldots<j_q}}\rho_{i_1}\ldots\rho_{i_p}\sigma_{j_1}\ldots\sigma_{j_q}\\
&=\frac{p!q!}{(p+q)!}\left(\sum_{i_1<\ldots<i_p}\rho_{i_1}\ldots\rho_{i_p}\right)
\left(\sum_{j_1<\ldots<j_q}\sigma_{j_1}\ldots\sigma_{j_q}\right)
=\frac{p!q!}{(p+q)!}\lVert\overset{p}{\Lambda}u\rVert_1\lVert\overset{q}{\Lambda}v\rVert_1 .
\end{aligned}
\]

\uline{Corollaire 1} Si $u,v\in L^1(E,F)$, alors \deleted{\unsure{soit}} pour $r\geq0$

\[
\deleted{\det}\det(1+\deleted{\unsure{r}}r|u+v|)\leq\det(1+r|u|)(\det(1+r|v|))\tag*{\textcircled{12}}
\]

i.e. on a

\[
\prod_i(1+r\rho_i(u+v))\leq\left(\prod_i1+r\rho_i(u)\right)\left(\prod_i(1+r\rho_i(v))\right)\tag*{\textcircled{13}}
\]

On va \unsure{même|en} montrer que l'on a majoration coëfficient par coëfficient, donc que

\uline{Corollaire 2} Pour tout $n$,

\[
\alpha_n(|u+v|)\leq\sum_{0\leq\underline{p}\leq n}\alpha_p(|u|)\alpha_{n-p}(|v|)\tag*{\textcircled{14}}
\]

En effet, on a \deleted{$|u+v|=U(u+v)$, où $U\in L(E,E)$, $\lVert U\rVert\leq1$, donc $\alpha_n(|u+v|)=\lVert\overset{n}{\Lambda}U(u+v)\rVert_1=\overset{n}{\Lambda}U\,\overset{n}{\Lambda}(u+v)$} \inserted{$\alpha_n(|u+v|)=\lVert\overset{n}{\Lambda}(u+v)\rVert_1$}

\[
\begin{aligned}
\overset{n}{\Lambda}(u+v)&=\sum_{p=0}^n\frac{n!}{p!(n-p)!}(\overset{p}{\Lambda}u)\wedge(\overset{n-p}{\Lambda}v)\qquad\text{d'où}\\
\alpha_n(|u+v|)\ \lVert\overset{n}{\Lambda}(u+v)\rVert_1
&\leq\sum_{p=0}^n\frac{n!}{p!(n-p)!}\lVert(\overset{p}{\Lambda}u)\wedge(\overset{n-p}{\Lambda}v)\rVert_1\leq\sum\deleted{\frac{n!}{\unsure{p!(n-p)}!}}\deleted{\unsure{p!q|\frac{p!(n-p)!}{n!}}}\alpha_{n}(|u|)\alpha_{n-p}(|v|).
\end{aligned}
\]

Prouv\unsure{ons|er} le

\uline{Théorème 2} Pour $u,v\in L_0(E,F)$, \inserted{$r\geq0$} et pour \deleted{tout} $m$, on a

\[
\prod_1^m(1+r\rho_i(u+v))\leq\left(\prod_1^m(1+r\rho_i(u))\right)\prod_1^m(1+r\rho_i(v))\tag*{\textcircled{15}}
\]

De façon plus précise, on a majoration coëfficients par coëfficients, i.e. pour \deleted{tout} entier \deleted{\unsure{$m$}} $m$, $0\leq m\leq n$, on a

\[
\deleted{S_n^m(u+v)\unsure{=|\leq}\sum_{p\unsure{=|\leq}0}^{m}}
\]

\[
\sigma_{\deleted{\unsure{p}}n}^{m}(u+v)\leq\sum_{p=0}^m\sigma_n^p(u)\,\sigma_{\text{--}}^{m-p}(v)\tag*{\textcircled{16}}
\]

% slanted
\margin{\uline{notation \unsure{?|$\leq0$}} à \uline{réviser}}

où de façon générale on pose, pour \inserted{$w\in L_0(E,F)$ et} $0\leq m\leq n$ :

\[
\sigma_n^m(w)=\sum_{1\leq i_1<\ldots<i_m\leq n}\rho_{i_1}(w)\ldots\rho_{i_m}(w)\tag*{\textcircled{17}}
\]

Comme \unsure{cas|en} particulier \unsure{des|aux} formules (16), pour $v=1$, signalons \inserted{\textcircled{18}}
